\bookchapter{נספח H: מפת הדרכים של הגרסאות}

\chapterintro{הספר הזה עדיין מתעדכן באופן שוטף. להלן הכיוונים המתוכננים לגרסאות העתידיות, ואנחנו מזמינים אותך לספר לנו במשוב מה התוכן שהכי תרצה לראות.}

\textbf{v1.8 כבר יצא לפועל: נוספו המדורים השביעי, השמיני והתשיעי — 'פיתוח מעשי בעזרת AI', 'פיתוח מעשי ל-Android' ו'פיתוח מעשי ל-iOS' (פרקים 37–69, 33 פרקים בסך הכול)}\par

מהמהפכה הפרדיגמטית של תכנות בעזרת AI, חמשת עקרונות החשיבה, הנדסת הקשר והנדסת דרישות, התמונה המלאה של הכלים והמודלים של 2026, עבודה מעשית עם API, הנדסת שיחה, שלושה פרויקטים שלמים, הבטחת איכות ודרך ההתקדמות — הספר יקפיץ אותך מ'משתמש ב-AI' ל'מנהל את ה-AI שכותב קוד'; וכן את שרשרת הפיתוח המלאה בשתי הפלטפורמות הגדולות, Android ו-iOS: הקמת סביבת עבודה (Android Studio / Xcode), מיני-קורס בשפות (Kotlin / Swift), פריימוורקים לממשק משתמש (Jetpack Compose / SwiftUI), נתונים ורשת, פרויקטים מעשיים (רישום הוצאות + בדיקת מזג אוויר), בדיקות ועלייה לחנויות (Google Play / App Store) ומדריכי עזר מהירים — כדי לקחת אותך ממשתמש דיגיטלי ועד למפתח אפליקציות.

\begin{tcbinfo}{הכיוונים המתוכננים לגרסה הבאה (v1.9)}
\begin{itemize}
  \item טלפון ומכשירים לבישים: מדריך לשיתוף פעולה מעמיק בין Apple Watch, שעוני Huawei וטבעות חכמות (Oura/Galaxy Ring)
  \item טיולים וטלפון: eSIM בפועל, פתרונות גלישה בינלאומיים, מפות אופליין, תשלומים בחו״ל והשוואת אפליקציות תרגום
  \item טלפון וניהול בריאות: כושר ובריאות, ניטור שינה, דופק וחמצן בדם, בריאות האישה ואפליקציות לבריאות הנפש
  \item יצירת סרטונים קצרים עם AI: מדריכים מעשיים ל-CapCut AI, Kling AI, אווטארים דיגיטליים ודיבוב בעזרת AI
  \item מבוא לדרמות קצרות בפורמט אנכי: מבנה תסריט, טכניקות צילום ומדיניות הפלטפורמות
\end{itemize}

\end{tcbinfo}

\begin{tcbwarn}{הערה על עדכניות התוכן}
טכנולוגיית הטלפונים מתקדמת בקצב מסחרר, ותוכן הספר נכתב בהתבסס על המצב הטכנולוגי והשוק באמצע 2026. יכולות ה-AI, דגמי השבבים, גרסאות המערכת וטווחי המחירים עשויים כבר להיות שונים ברגע שתקרא. בכל הנוגע לצעדים מעשיים, סמכו על גרסת המערכת בפועל של המכשיר שלכם; בכל הנוגע להחלטות רכישה, בדקו ביקורות ומשובי משתמשים עדכניים.

\end{tcbwarn}
